\scanpage{11}
\mspage{4}

Le groupe $\mathbb Z/6\mathbb Z$ opère transitivement sur $({}_2X_S)\setminus\{0\}$ car $\pm1$ opère trivialt - donc \({}_2X_S\setminus\{0\}\) est un torseur sous
\[
(\mathbb Z/6\mathbb Z)/\pm1\simeq\mathbb Z/3\mathbb Z
\]
\deleted{(\unsure{conj.}} (générateur canonique $\unsure{\simeq}\,-\bar\zeta\,\unsure{j}=\exp\frac{2\pi}{6}$),
\deleted{\unsure{en ce cas}} \inserted{le torseur} s'identifie \unsure{:|à} $S/\pm1$, en associant à chaque $s\in S$ le milieu du rayon $[0,s]$, ou plutôt son image dans $X_S$.

L'ens des pts fixes de $\mathbb Z/3\mathbb Z\inserted{\subset\mathbb Z/6\mathbb Z}$ sur $X_S$ est de card 3, il s'obtient en prenant les barycentres des six triangles de la subdivision centrale (\unsure{bibarycentrique}) de l'hexagone (pts cerclés) (tels que $(1-\bar\zeta)/3$ etc dans le cas épinglé), ils se groupent en deux paquets de 3 \unsure{mutuellement} congrus \unsure{sous|mod} $\Gamma$, \unsure{donnant|donnent} deux éléments (symétriques) de $X_S$, d'ordre 3,

\margin{10}
